\section{Implementación Reproducible en Python}

% ================================================================
\begin{frame}[fragile]{Pipeline con Scikit-Learn: Prevención de Data Leakage}
\begin{columns}[T,onlytextwidth]
\begin{column}{0.58\textwidth}
\begin{lstlisting}[style=ccdepython]
from sklearn.datasets import load_breast_cancer
from sklearn.model_selection import train_test_split
from sklearn.pipeline import Pipeline
from sklearn.preprocessing import StandardScaler
from sklearn.neighbors import KNeighborsClassifier

X, y = load_breast_cancer(return_X_y=True)
X_tr, X_te, y_tr, y_te = train_test_split(
    X, y, test_size=0.25, stratify=y, random_state=42
)

# Pipeline estricto: escalamiento + estimador
pipe = Pipeline([
    ("scaler", StandardScaler()),
    ("knn", KNeighborsClassifier())
])
\end{lstlisting}
\end{column}

\begin{column}{0.40\textwidth}
\begin{keybox}
\textbf{Arquitectura Pipeline}:
\begin{itemize}\setlength{\itemsep}{0.15em}
  \footnotesize
  \item Aplica $\mu$ y $\sigma$ aprendidos \textbf{solo en el pliegue de entrenamiento}.
  \item Evita fuga de información (\textit{data leakage}) al test/validación.
  \item Facilita la reproducibilidad integral del flujo.
\end{itemize}
\end{keybox}

\vspace{0.10cm}
\begin{ccdebox}
\scriptsize
\textbf{Dataset}: Breast Cancer Wisconsin (569 obs., 30 features numéricas).
\end{ccdebox}
\end{column}
\end{columns}
\source{Script reproducible \texttt{codigo\_demo\_knn.py} (CCDE).}
\end{frame}

% ================================================================
\begin{frame}[fragile]{Búsqueda de Hiperparámetros con GridSearchCV}
\begin{columns}[T,onlytextwidth]
\begin{column}{0.58\textwidth}
\begin{lstlisting}[style=ccdepython]
from sklearn.model_selection import GridSearchCV

param_grid = {
    "knn__n_neighbors": [3, 5, 7, 9, 11, 15],
    "knn__weights": ["uniform", "distance"],
    "knn__metric": ["minkowski"],
    "knn__p": [1, 2] # 1: Manhattan, 2: Euclidiana
}

grid = GridSearchCV(
    pipe,
    param_grid=param_grid,
    cv=5,
    scoring="roc_auc",
    n_jobs=-1
)
grid.fit(X_tr, y_tr)
best_knn = grid.best_estimator_
\end{lstlisting}
\end{column}

\begin{column}{0.40\textwidth}
\begin{contrastbox}{ccdemid}{Espacio de Búsqueda}
\footnotesize
\begin{itemize}\setlength{\itemsep}{0.15em}
  \item \textbf{Vecindades ($k$)}: de $3$ a $15$ impares para evitar empates.
  \item \textbf{Ponderación}: uniforme vs inversa a la distancia.
  \item \textbf{Métricas}: norma $L_1$ vs $L_2$.
  \item \textbf{Criterio}: ROC-AUC con 5 pliegues estratificados.
\end{itemize}
\end{contrastbox}
\end{column}
\end{columns}
\source{Pedregosa et al. (2011); scikit-learn documentation.}
\end{frame}

% ================================================================
\begin{frame}{Resultados Empíricos y Matriz de Confusión}
\begin{columns}[T,onlytextwidth]
\begin{column}{0.51\textwidth}
\begin{ccdebox}
\footnotesize
\textbf{Mejor Configuración Encontrada}
\vspace{0.06cm}

\renewcommand{\arraystretch}{1.12}
\begin{tabular*}{\linewidth}{@{\extracolsep{\fill}}l r@{}}
\toprule
\textbf{\textcolor{ccdenavy}{Parámetro / Métrica}} & \textbf{\textcolor{ccdenavy}{Valor}} \\
\midrule
$k$ óptimo (\texttt{n\_neighbors}) & \textbf{15} \\
Ponderación (\texttt{weights}) & \textbf{\texttt{distance}} \\
Métrica & \textbf{Euclidiana ($L_2$)} \\
ROC-AUC CV (5-fold) & \textbf{0.9875} \\
Accuracy en Test & \textbf{96.50\%} \\
ROC-AUC en Test & \textbf{0.9948} \\
\bottomrule
\end{tabular*}
\end{ccdebox}
\end{column}

\begin{column}{0.47\textwidth}
\begin{contrastbox}{okgreen}{Matriz de Confusión en Test}
\footnotesize
\vspace{0.04cm}

\renewcommand{\arraystretch}{1.12}
\begin{tabular*}{\linewidth}{@{\extracolsep{\fill}}r c c@{}}
\toprule
& \textbf{Pred. Mal.} & \textbf{Pred. Ben.} \\
\midrule
\textbf{Real Maligno} & \textbf{48} & 5 \\
\textbf{Real Benigno} & 0 & \textbf{90} \\
\bottomrule
\end{tabular*}

\vspace{0.06cm}
\begin{itemize}\setlength{\itemsep}{0.08em}
  \scriptsize
  \item Sensibilidad (Recall benigno): \textbf{100\%}.
  \item Especificidad (clase maligna): \textbf{90.57\%}.
\end{itemize}
\end{contrastbox}

\vspace{0.06cm}
\begin{keybox}
\scriptsize \textbf{Conclusión}: La ponderación por distancia combinada con $k=15$ maximizó la generalización en test.
\end{keybox}
\end{column}
\end{columns}
\source{Resultados reproducibles obtenidos con \texttt{codigo\_demo\_knn.py} (semilla 42).}
\end{frame}
